% p028_d 的电脑重画（与 p028_c 等效的画法：①②③串联接电池，④并联在“②+③”两端，电容并联③）
\begin{tikzpicture}[line width=0.7pt, >={Stealth[length=2.2mm]}]
  % 主回路：电池在底边
  \draw (0,0) -- (3.45,0);
  \draw (3.45,-0.4) -- (3.45,0.4);
  \draw (3.7,-0.2) -- (3.7,0.2);
  \draw (3.7,0) -- (7.4,0) -- (7.4,2.9);
  \draw (0,0) -- (0,1.3) -- (1.0,1.3);
  \draw (1.0,1.12) rectangle (2.1,1.48);
  \draw (2.1,1.3) -- (3.1,1.3);
  \draw (3.1,1.12) rectangle (4.2,1.48);
  \draw (4.2,1.3) -- (5.5,1.3);
  \draw (5.5,1.12) rectangle (6.6,1.48);
  \draw (6.6,1.3) -- (7.4,1.3);
  % ④：并联在②+③两端
  \draw (2.55,1.3) -- (2.55,2.9) -- (4.0,2.9);
  \draw (4.0,2.72) rectangle (5.1,3.08);
  \draw (5.1,2.9) -- (7.4,2.9);
  % 电容：并联③
  \draw (4.9,1.3) -- (4.9,0.55) -- (6.15,0.55);
  \draw (6.15,0.2) -- (6.15,0.9);
  \draw (6.4,0.2) -- (6.4,0.9);
  \draw (6.4,0.55) -- (7.4,0.55);
  % 结点
  \fill (2.55,1.3) circle (1.6pt) (4.9,1.3) circle (1.6pt)
        (7.4,1.3) circle (1.6pt) (7.4,0.55) circle (1.6pt);
  % 编号
  \node at (1.55,1.95) [draw, circle, inner sep=1pt] {1};
  \node at (3.65,1.95) [draw, circle, inner sep=1pt] {2};
  \node at (6.05,1.95) [draw, circle, inner sep=1pt] {3};
  \node at (4.55,2.42) [draw, circle, inner sep=1pt] {4};
\end{tikzpicture}
